\documentclass[11pt,a4paper]{article}

%% PACHETE MATEMATICE
\usepackage{amsmath,amssymb,amsfonts}
\usepackage{amsthm}
\usepackage{mathpazo}
\usepackage{eulervm}
\usepackage{enumerate}
\usepackage{color}
\usepackage{graphicx}
\usepackage[all]{xy}
\usepackage{xcolor}
\usepackage{framed}
\usepackage[unicode]{hyperref}
\setcounter{MaxMatrixCols}{20}
\usepackage{multicol}
%\usepackage[romanian]{babel}


%% ASPECT
\linespread{1.05}
\usepackage[T1]{fontenc}
\usepackage[utf8x]{inputenc}
\usepackage[margin=2cm]{geometry}
% \usepackage[color]{showkeys}
\usepackage{anyfontsize}
\usepackage{titlesec}
\titleformat*{\section}{\large\bfseries}

\usepackage{fancyhdr}
\pagestyle{fancy}
\rhead{\small \color{gray} Notițe de Adrian Manea}
\lhead{\small \color{gray} Aero-BCD, 2017---2018, Prof. L. Costache \& M. Olteanu}


\usepackage[autostyle = false, style = german]{csquotes}
\usepackage[romanian]{babel}
\newcommand\qq{\enquote}

%% COMENZILE MELE
\renewcommand*{\proofname}{Dem.:}
\newcommand\bb{\mathbb}
\newcommand\dr{\mathrm}
\newcommand\kal{\mathcal}
\newcommand\fk{\mathfrak}
\newcommand\op{\oplus}
\newcommand\ot{\otimes}
\newcommand\ds{\displaystyle}
\newcommand\id{\indent}
\newcommand\nid{\noindent}
\newcommand\seq{\subseteq}
\newcommand\sneq{\subsetneq}
\newcommand\speq{\supseteq}
\newcommand\spneq{\supsetneq}
\newcommand\rar{\rightarrow}
\newcommand\Rar{\Rightarrow}
\newcommand\xrar{\xrightarrow}
\newcommand\lar{\leftarrow}
\newcommand\Lar{\Leftarrow}
\newcommand\xlar{\xleftarrow}
\newcommand\lrar{\leftrightarrow}
\newcommand\Lrar{\Leftrightarrow}
\newcommand\ideal{\trianglelefteq}
\newcommand\ai{\text{ a.\^{i}. }}
\newcommand\sk{(\textit{Schi\c{t}\u{a}}) }
\newcommand\rhup{\rightharpoonup}
\newcommand\rhdn{\rightharpoondown}
\newcommand\lhup{\leftharpoonup}
\newcommand\lhdn{\leftharpoondown}
\newcommand\vid{\emptyset}
\newcommand\wh{\widehat}
\newcommand\ol{\overline}
\newcommand\NN{\mathbb{N}}
\newcommand\RR{\mathbb{R}}
\newcommand\ZZ{\mathbb{Z}}
\newcommand\QQ{\mathbb{Q}}
\newcommand\CC{\mathbb{C}}

%% STILURILE MELE
\newtheoremstyle{dotless}{}{}{\itshape}{}{\bfseries}{:}{ }{}

\theoremstyle{dotless}
\newtheorem{theorem}{Teorem\u{a}}[section]
\newtheorem{proposition}{Propozi\c{t}ie}[section]
\newtheorem{lemma}{Lem\u{a}}[section]
\newtheorem{corollary}{Corolar}[section]
\newtheorem{exercise}{Exerci\c{t}iu}

\newtheoremstyle{dotlessdr}{}{}{}{}{\bfseries}{:}{ }{}
\theoremstyle{dotlessdr}
\newtheorem{example}{Exemplu}[section]
\newtheorem{definition}{Defini\c{t}ie}[section]
\newtheorem{remark}{Observa\c{t}ie}[section]




\begin{document}
% \definecolor{labelkey}{RGB}{222,0,0}

\begin{center}
  {\large\textbf{Seminar 5 \\ Șiruri și serii de funcții. Serii de puteri}}
\end{center}

\section{Șiruri de funcții}
\label{sec:siruri-f}

\begin{definition}\label{def:pc-uc}
  Fie $(f_n)_n$ un șir de funcții, cu fiecare $f_n : [a, b] \to \RR$ și fie
  o funcție $f: [a, b] \to \RR$. Spunem că șirul $(f_n)$ \emph{converge punctual}
  pe $[a, b]$ către $f$ pentru $n \to \infty$, scris $f_n \xrar{PC} f$ dacă avem
  $f_n(x_0) \to f(x_0)$, pentru orice $x \in [a, b]$.

  Spunem că șirul $(f_n)$ \emph{converge uniform} pe $[a, b]$ către $f$, scris
  $f_n \xrar{UC} f$, dacă:
  \[
    \forall \varepsilon > 0, \exists N_\varepsilon \in \NN \ai %
    \forall n \geq N_\varepsilon, |f_n(x) - f(x)| < \varepsilon, \forall x \in [a, b].
  \]
\end{definition}

În calcule, va fi de folos următorul rezultat:

\begin{theorem}\label{thm:pc-uc}
  Un șir $(f_n)$ de funcții mărginite $f_n : [a, b] \to \RR$ este uniform convergent
  către o funcție $f$ dacă și numai dacă $\ds\lim_{n \to \infty} ||f_n - f|| = 0$.
\end{theorem}

De asemenea, avem și:
\begin{theorem}\label{thm:pc-uc-implica}
  Orice șir de funcții $f_n : [a, b] \to \RR$ uniform convergent pe $[a, b]$ este
  punctual convergent pe $[a, b]$.
\end{theorem}

Reciproca acestui rezultat este falsă: fie $[a, b] = [0, 1]$ și $f_n(x) = x^n, n \geq 1$.

Pentru orice $x \in [a, b]$, avem:
\[
  \lim_{n \to \infty} f_n(x) =
  \begin{cases}
    0, & x \in [0, 1) \\
    1, & x = 1
  \end{cases}
\]
Rezultă că $f_n \xrar{PC} f$, unde funcția $f$ este definită de limita de mai sus.
Dar:
\begin{align*}
  ||f_n - f|| &= \sup_{x \in [0, 1)} |f_n(x) - f(x)| \\
              &= \max (\sup_{x \in [0, 1)} (|f_n(x) - f(x)|, |f_n(1) - f(1)|)) \\
              &= \max (\sup_{x \in [0, 1)} (x^n, 0)) \\
              &= 1,
\end{align*}
de unde obținem $\ds\lim_{n \to \infty} ||f_n - f|| = 1 \neq 0$. Rezultă că
$(f_n)$ este punctual convergent, nu uniform convergent pe $[0, 1)$.

\vspace{1cm}

Pentru caracterizarea convergenței, avem la dispoziție mai multe rezultate,
printre care și \emph{criteriul fundamental, al lui Cauchy}. Vom enunța, însă,
rezultatele care vor fi utile în special în rezolvarea exercițiilor:

\begin{theorem}[Integrare termen cu termen]\label{thm:uc-integrala}
  Fie $(f_n)$ un șir uniform convergent de funcții continue, $f_n : [a, b] \to \RR$.

  Atunci limita $f = \ds\lim_{n \to \infty} f_n$ este o funcție continuă pe
  $[a, b]$ și avem:
  \[
    \lim_{n \to \infty} \int_a^b f_n(x) dx = \int_a^b f(x) dx.
  \]

  Cu alte cuvinte, în cazul convergenței uniforme, limita comută cu integrala.
\end{theorem}

Pentru derivabilitate, avem:

\begin{theorem}[Derivare termen cu termen]\label{thm:uc-derivata}
  Fie $(f_n)$ un șir de funcții din $C^1([a, b])$ și $f, g$ funcții mărginite,
  cu $f, g : [a, b] \to \RR$. Dacă $f_n \xrar{PC} f$ și $f'_n \xrar{UC} g$ pe
  $[a, b]$, atunci $f$ este derivabilă pe $[a, b]$ și $f' = g$.
\end{theorem}

Pentru cazul șirurilor monotone de funcții continue, rezultatul următor arată
o legătură simplă între convergența punctuală și cea uniformă:

\begin{theorem}[U. Dini]\label{thm:dini}
  Fie $(f_n)$ un șir monoton de funcții continue, $f_n : [a, b] \to \RR$, astfel
  încît $f_n \xrar{PC} f$. Atunci $f_n \xrar{UC} f$.
\end{theorem}

Următorul rezultat va fi de folos în special în capitolul privitor la serii
de puteri:
\begin{theorem}[Stone---Weierstrass]\label{thm:label}
  Pentru orice funcție continuă $f : [a, b] \to \RR$, există un șir $(f_n)$
  \emph{de polinoame}, cu $f_n : [a, b] \to \RR$, astfel încît $f_n \xrar{UC} f$.
\end{theorem}

%%%%%%%%%%%%%%%%%%%%%%%%%%%%%%%%%%%%%%%%%%%%%%%%%%%%%%%%%%%%%%%%%%%%%%

\section{Serii de funcții}
\label{sec:serii-f}

Trecem acum la studiul seriilor de funcții, făcînd legături similare cu trecerea
de la șiruri de numere la serii de numere.

Începem cu o noțiune fundamentală:

\begin{definition}\label{def:pc-uc-serie-f}
  Mulțimea valorilor lui $x$ pentru care seria de funcții $\ds\sum_{n \geq 1} f_n(x)$
  este convergentă se numește \emph{mulțimea de convergență} a seriei, iar funcția
  $f : [a, b] \to \RR$, cu $f(x) = \ds\lim_{n \to \infty} S_n(x)$, unde $S_n(x)$
  este șirul sumelor parțiale pentru seria de funcții, se numește \emph{suma seriei}.

  
  Seria $\sum f_n$ se numește \emph{simplu (punctual) convergentă} către funcția $f$
  dacă șirul sumelor parțiale $(S_n(x))_n$ este punctual convergent către $f$.

  Similar, seria se numește \emph{uniform convergentă} către funcția $f$ dacă șirul
  sumelor parțiale este uniform convergent către $f$.

  Seria se numește \emph{absolut convergentă} dacă seria $\sum |f_n(x)|$ este simplu
  convergentă.
\end{definition}

Regăsim acum, atît noțiunile privitoare la șirurile de funcții, cît și criteriile
de convergență pentru serii de numere.

\begin{theorem}\label{thm:serie-f-integrala}
  Fie $\sum f_n$ o serie uniform convergentă de funcții continue, cu $f_n : [a, b] \to \RR$
  și fie $s$ suma acestei serii. Atunci $s$ este o funcție continuă pe $[a, b]$ și:
  \[
    \int_a^b s(x) dx = \sum_{n \geq 1} \int_a^b f_n(x) dx.
  \]
\end{theorem}

Avem rezultatul corespunzător și pentru derivate:

\begin{theorem}\label{thm:serie-f-derivata}
  Fie $\sum f_n$ o serie punctual convergentă de funcții din $C^1([a, b])$, cu suma $s$
  pe $[a, b]$ și astfel încît seria derivatelor $\sum f'_n$ să fie uniform convergentă.
  Atunci funcția $s$ este derivabilă pe $[a, b]$ și în plus:
  \[
    s'(x) = \sum_{n \geq 1} f'_n(x).
  \]
\end{theorem}

\begin{theorem}[Weierstrass]\label{thm:weierstrass-serii-f}
  Fie $\sum f_n$ o serie de funcții, cu $f_n : [a, b] \to \RR$ și fie $\sum a_n$ o serie
  convergentă de numere reale pozitive.

  Dacă $|f_n(x)| \leq a_n$, pentru orice $x \in [a, b]$ și $n \geq N$, cu $N$ fixat,
  atunci seria de funcții $\sum f_n$ este uniform convergentă pe $[a, b]$.
\end{theorem}

\begin{theorem}[Criteriul lui Abel]\label{thm:abel-serii-f}
  Dacă seria de funcții $\sum f_n$ se poate scrie $\sum \alpha_n v_n$, astfel încît
  seria de funcții $\sum v_n$ să fie uniform convergentă, iar $(\alpha_n)$ să fie un
  șir monoton de funcții egal mărginite (i.e.\ mărginite de aceeași constantă),
  atunci seria inițială este uniform convergentă.
\end{theorem}

O alternativă:

\begin{theorem}[Criteriul lui Dirichlet]\label{thm:dirichlet-serii-f}
  Dacă seria de funcții $\sum f_n$ se poate scrie sub forma $\sum \alpha_n v_n$,
  astfel încît șirul sumelor parțiale al seriei $\sum v_n$ să fie un șir de funcții
  egal mărginite, iar șirul $(\alpha_n)$ să fie un șir monoton de funcții, ce converge
  uniform către 0, atunci ea este uniform convergentă.
\end{theorem}

%%%%%%%%%%%%%%%%%%%%%%%%%%%%%%%%%%%%%%%%%%%%%%%%%%%%%%%%%%%%%%%%%%%%%%

\section{Formula lui Taylor}
\label{sec:taylor}

Putem asocia oricărei funcții cu anumite proprietăți \qq{bune} un polinom care o
aproximează. Este vorba despre \emph{polinomul Taylor}, definit astfel.

Fie $I \seq \RR$ un interval deschis și fie $f : I \to \RR$ o funcție de clasă
$C^m(I)$. Pentru orice $a \in I$, definim \emph{polinomul Taylor} de gradul $n \leq m$
asociat funcției $f$ în punctul $a$:
\[
  T_{n, f, a} = \sum_{k = 0}^n \frac{f^{(k)}(a)}{k!} (x - a)^k.
\]

Restul (eroarea aproximării) se definește prin:
\[
  R_{n, f, a}(x) = f(x) - T_{n, f, a} (x).
\]

Următorul rezultat ne arată că polinomul de mai sus poate fi transformat într-o formulă
mai exactă:

\begin{theorem}[Formula lui Taylor cu resturi Lagrange]\label{thm:taylor-lagrange}
  Fie $f : I \to \RR$ o funcție de clasă $C^{n+1}(I)$ și $a \in I$. Atunci, pentru
  orice $x \in I$, există $\xi \in (a, x)$ (sau $(x, a)$, după caz), astfel încît:
  \[
    f(x) = T_{n, f, a}(x) + \frac{(x - a)^{n+1}}{(n + 1)!} f^{n + 1}(\xi).
  \]
\end{theorem}

Privitor la restul rezultat din această formulă, avem următoarele:
\begin{itemize}
\item \emph{forma Peano}: $\exists \omega : I \to \RR$, cu
  $\ds\lim_{x \to a} \omega(x) = \omega(a) = 0$:
  \[
    R_{n, f, a}(x) = \frac{(x - a)^n}{n!} \omega(x).
  \]
\item \emph{forma integrală}:
  \[
    R_{n, f, a}(x) = \frac{1}{n!} \int_a^x f^{(n + 1)}(t) (x - t)^n dt.
  \]
\item $\ds\lim_{x \to a} \frac{R_{n, f, a}(x)}{(x - a)^n} = 0$.
\end{itemize}

%%%%%%%%%%%%%%%%%%%%%%%%%%%%%%%%%%%%%%%%%%%%%%%%%%%%%%%%%%%%%%%%%%%%%%

\section{Exerciții șiruri și serii de funcții}
\label{sec:exercitii}

1. Să se studieze convergența simplă și uniformă a șirurilor de funcții:
\begin{enumerate}[(a)]
\item $f_n : [-1, 1] \to \RR, f_n(x) = \frac{x}{1 + nx^2}$;
\item $f_n : [0, 1] \to \RR, f_n(x) = x^n(1 - x^n)$;
\item $f_n : (-1, 1) \to \RR, f_n(x) = \frac{1 - x^n}{1 - x}$;
\item $f_n : [0, 1] \to \RR, f_n(x) = \frac{2nx}{1 + n^2 x^2}$;
\item $f_n : [0, \infty) \to \RR, f_n(x) = \frac{x + n}{x + n + 1}$;
\item $f_n : [0, \infty) \to \RR, f_n(x) = \frac{x}{1 + nx^2}$;
\item $f_n : \RR \to \RR, f_n(x) = \arctan (nx)$;
\item $f_n : \big[ 0, \frac{\pi}{2} \big] \to \RR, f_n(x) = n \sin^n x \cos x$.
\end{enumerate}

\vspace{1cm}

2. Să se arate că șirul $f_n : [0, \infty) \to \RR, f_n(x) = \frac{x}{x + n}$ nu converge
uniform pe $[0, \infty)$, dar converge uniform pe orice interval $[a, b]$, cu $0 < a < b$.

\vspace{1cm}

3. Să se arate că șirul de funcții $(f_n)$, unde:
\[
  f_n : (1, \infty) \to \RR, \quad f_n(x) = \sqrt{(n^2 + 1) \sin^2 \frac{\pi}{n} + nx} - \sqrt{nx}
\]
este uniform convergent (la 0).

\vspace{1cm}
4. Să se arate că șirul de funcții $(f_n)$, cu:
\[
  f_n : \RR \to \RR, \quad f_n(x) = \frac{1}{n} \arctan x^n
\]
converge uniform pe $\RR$, dar:
\[
  \Big( \lim_{n \to \infty} f_n(x) \Big)'_{x = 1} \neq \lim_{n \to \infty} f'_n(1).
\]

Rezultatele diferă deoarece șirul derivatelor nu converge uniform pe $\RR$.

\vspace{1cm}

5. Să se arate că șirul de funcții $(f_n)$, cu:
\[
  f_n : [0, 1] \to \RR, \quad f_n(x) = nxe^{-nx^2}
\]
converge, dar:
\[
  \lim_{n \to \infty} \int_0^1 f_n(x) dx \neq \int_0^1 \lim_{n \to \infty} f_n(x) dx.
\]

Rezultatul se explică prin faptul că șirul nu este uniform convergent. De exemplu,
pentru $x_n = \frac{1}{n} \in [0, 1]$, avem $f_n(x_n) \to 1$, dar în general,
$f_n(x) \to 0$.

\vspace{1cm}

6. Studiați convergența simplă și uniformă a șirului de funcții:
\[
  f_n : \Big[ 0, \frac{\pi}{2} \Big] \to \RR, \quad %
  f_n(x) = n \sin^n x \cos x.
\]

%%%%%%%%%%%%%%%%%%%%%%%%%%%%%%%%%%%%%%%%%%%%%%%%%%%%%%%%%%%%%%%%%%%%%%
\section{Serii de puteri. Serii Taylor}
\label{sec:taylor}

6. Să se dezvolte în seria Maclaurin următoarele funcții, precizînd și
domeniul de convergență:
\begin{enumerate}[(a)]
\item $f(x) = e^x$;
\item $f(x) = \sin x$.
\end{enumerate}

\emph{Soluție:} (a) Calculăm derivatele funcției și obținem că $(e^x)^{(n)} = e^x$,
pentru orice $n$. Atunci, din formula lui Maclaurin, găsim:
\[
  e^x = \sum_{n \geq 0} \frac{1}{n!} \cdot x^n.
\]
Pentru a determina domeniul de convergență, folosim criteriul raportului:
\[
  \lim_{n \to \infty} \Big| \frac{x^{n+1}}{(n+1)!} \cdot \frac{n!}{x^n} \Big| = %
  \lim_{n \to \infty} \Big| \frac{x}{n+1} \Big| = 0 < 1,
\]
deci seria este convergentă pe $\RR$.

Restul de ordin $n$ se obține:
\[
  R_n(x) = \frac{x^{n + 1}}{(n + 1)!} \cdot e^{\xi}, \quad \xi \in (0, x) \text{ sau } %
  \xi \in (x, 0).
\]

(b) Calculăm derivata de ordin $n$ a funcției sinus și observăm:
\[
  (\sin x)^{(n)} = \sin \Big( x + n \frac{\pi}{2} \Big), \forall n \in \NN.
\]
Deducem că derivatele de ordin par sînt nule, iar cele de ordin impar sînt $(-1)^n$.
Rezultă:
\[
  \sin x = \sum_{n \geq 0} \frac{(-1)^n}{(2n + 1)!} \cdot x^{2n + 1}.
\]
Pentru determinarea domeniului de convergență, putem folosi din nou criteriul
raportului și obținem:
\[
  \lim_{n \to \infty} \Big| \frac{x^{2n + 3}}{(2n + 3)!} \cdot \frac{(2n + 1)!}{x^{2n + 1}} \Big| = %
  \lim_{n \to \infty} \frac{x^2}{(2n + 2)(2n + 3)} = 0 < 1,
\]
deci seria este convergentă pe $\RR$.

\vspace{1cm}

7. Să se dezvolte în serie de puteri ale lui $x - 1$ funcția:
\[
  f : \RR - \{ 0 \} \to \RR, \quad f(x) = \frac{1}{x}.
\]

\emph{Soluție:} Derivata de ordin $n$ este:
\[
  f^{(n)}(x) = \frac{(-1)^n n!}{x^{n+1}} \Rightarrow f^{(n)}(1) = (-1)^n n!, \forall n \in \NN.
\]
Rezultă:
\[
  \frac{1}{x} = \sum_{n \geq 0} (-1)^n (x - 1)^n, \forall x \in (0, 2).
\]

\vspace{1cm}

8. Să se afle raza de convergență și mulțimea de convergență pentru seriile de puteri:
\begin{enumerate}[(a)]
\item $\ds\sum_{n \geq 0} x^n$;
\item $\ds\sum_{n \geq 1} \frac{n^n x^n}{n!}$;
\end{enumerate}

\emph{Soluție:} (a) Fie $R$ raza de convergență. Din definiție, avem:
\[
  R = \lim_{n \to \infty} \frac{|a_n|}{|a_{n + 1}|} = 1.
\]
Deducem că seria este absolut convergentă pe $(-1, 1)$ și divergentă în rest.

Evident, seria este uniform convergentă pe orice interval închis $|x| \leq r < 1$, iar
pentru $x = \pm 1$, obținem o serie divergentă.

(b) Din nou, calculăm raza de convergență:
\[
  R = \lim_{n \to \infty} \frac{|a_n|}{|a_{n + 1}|} = %
  \lim_{n \to \infty} \Big( \frac{n}{n + 1} \Big)^n = \frac{1}{e}.
\]
Rezultă că seria este absolut convergentă pe $\Big( -\frac{1}{e}, \frac{1}{e} \Big)$
și divergentă în rest. Seria este uniform convergentă pe $|x| \leq r < e$, iar
pentru $x = \pm \frac{1}{e}$, se obține o serie divergentă, rezultat care
se poate demonstra cu criteriul raportului.

\vspace{1cm}

9. Studiați convergența seriilor de funcții:
\begin{enumerate}[(a)]
\item $\ds\sum_{n \geq 0} \ln^n x, x > 0$;
\item $\ds\sum_{n \geq 1} \Big( 1 + \frac{1}{n} \Big)^{-n^2} e^{-nx}$;
\item $\ds\sum_{n \geq 1} \frac{1}{n! x^n}$.
\end{enumerate}

\emph{Soluție:} Exercițiile se rezolvă cu o schimbare de variabilă, care
transformă seriile de funcții în serii de puteri:
\begin{enumerate}[(a)]
\item $t = \ln x$;
\item $t = e^{-x}$;
\item $t = x^{-1}$.
\end{enumerate}

\vspace{1cm}

10. Să se calculeze sumele seriilor de funcții:
\begin{enumerate}[(a)]
\item $\ds\sum_{n \geq 0} x^{2n}$;
\item $\ds\sum_{n \geq 0} (-1)^n \frac{x^{2n + 1}}{2n + 1}$;
\item $\ds\sum_{n \geq 0} (n + 1)x^n$.
\end{enumerate}

\emph{Soluție:} (a) Putem folosi direct formula pentru seria geometrică:
\[
  \sum_{n \geq 0} x^{2n} = \sum_{n \geq 0} (x^2)^n = \frac{1}{1 - x^2}, %
  \forall x \in (-1, 1).
\]

(b) Raza de convergență se poate calcula și obținem $R = 1$. Fie $S(x)$ suma seriei
date. Seria derivatelor are aceeași rază de convergență cu seria inițială, iar
suma ei este egală cu derivata sumei (seriile de puteri se pot deriva mereu termen
cu termen):
\[
  S'(x) = \sum_{n \geq 0} (-1)^n (x^2)^n = \frac{1}{1 + x^2}, \forall x \in (-1, 1).
\]
Rezultă $S(x) = \arctan x + C$ și, deoarece $S(0) = 0$, avem $C = 0$.

(c) Procedăm similar, acum cu integrare termen cu termen. Raza de convergență
este egală cu 1. Avem:
\[
  \int_0^x S(x) dx = \sum_{n \geq 0} x^{n + 1} = \frac{x}{1 - x}, \forall x \in (-1, 1).
\]
Rezultă $S(x) = \frac{1}{(1 - x)^2}$.

%%%%%%%%%%%%%%%%%%%%%%%%%%%%%%%%%%%%%%%%%%%%%%%%%%%%%%%%%%%%%%%%%%%%%%

\section{Exerciții serii de puteri}
\label{sec:ex}

1. Să se studieze convergența punctuală și uniformă a șirurilor de funcții:
\begin{enumerate}[(a)]
\item $f_n : \RR \to \RR, f_n(x) = \frac{nx^2}{1 + nx^2}$;
\item $f_n : [0, 1] \to \RR, f_n(x) = \frac{x^n}{1 + x^{2n}}$;
\item $f_n : \RR_+ \to \RR, f_n(x) = e^{-nx} \sin nx$;
\end{enumerate}

\vspace{1cm}

2. Să se afle mulțimea de convergență a seriei de funcții:
\[
  \sum_{n \geq 1} (-1)^n \frac{1}{n^x}.
\]

\vspace{1cm}

3. Să se studieze convergența seriilor:
\begin{enumerate}[(a)]
\item $\ds\sum_{n \geq 1} \frac{1}{x^2 + n^2}$;
\item $\ds\sum_{n \geq 1} \frac{\sin nx}{\sqrt[3]{n^4 + x^4}}$;
\item $\ds\sum_{n \geq 1} \frac{1}{x^2 + 2^n}, x \in \RR$.
\end{enumerate}

\vspace{1cm}

4. Să se determine mulțimea de convergență și suma seriilor:
\begin{enumerate}[(a)]
\item $\ds \sum_{n \geq 1} (-1)^{n + 1} \frac{x^n}{n}$;
\item $\ds \sum_{n \geq 0} (n + 1)(n + 2)(n + 3) x^n$;
\item $\ds \sum_{n \geq 1} n^3 x^n$;
\end{enumerate}

\vspace{1cm}

5. Găsiți mulțimea de convergență pentru seriile de puteri:
\begin{enumerate}[(a)]
\item $\ds\sum_{n \geq 1} \frac{1}{n! x^n}$;
\item $\ds\sum_{n \geq 1} \frac{x^n}{2^n + 3^n}$;
\item $\ds\sum_{n \geq 1} \sqrt[n]{n!} x^n$;
\item $\ds\sum_{n \geq 1} \frac{1 \cdot 5 \cdot 9 \cdots (4n - 3)}{ %
    3 \cdot 7 \cdot 11 \cdots (4n - 1)} x^n$;
\end{enumerate}

\vspace{1cm}

6. Să se dezvolte în serie Maclaurin următoarele funcții, precizînd domeniul
de convergență:
\begin{enumerate}[(a)]
\item $f(x) = \cos x$;
\item $f(x) = (1 + x)^a, a \in \RR$;
\item $f(x) = \frac{1}{1 + x}$;
\item $f(x) = \sqrt{1 + x}$;
\item $f(x) = \ln ( 1 + x)$;
\item $f(x) = \arctan x$.
\end{enumerate}

\vspace{1cm}

7. Să se dezvolte în serie de puteri ale lui $x + 4$ funcția:
\[
  f: \RR - \{ -2, -1 \} \to \RR, \quad %
  f(x) = \frac{1}{x^2 + 3x + 2}.
\]




\end{document}